\subsection{注记与进一步结果}
\label{sec:ch4-notes-further-results}

\subsubsection{参考文献}
\label{subsec:ch4-references}

旋转方法由 A. P. Calderón 与 A. Zygmund 在论文 \emph{On singular integrals}, Amer. J. Math. 78 (1956), 289--309 中引入. 第~\ref{sec:ch4-definition-examples}~节的例子来自他们更早的论文 \emph{On the existence of certain singular integrals}, Acta Math. 88 (1952), 85--139, 该文将在\MBUnresolvedReference{chapter}{第 5 章}讨论. 偶核情形下定理~\ref{thm:ch4-rough-kernel-lp} 的证明取自 A. P. Calderón 的课程讲义 \emph{Singular integrals and their applications to hyperbolic differential equations}, University of Buenos Aires, 1960. Zygmund \MBUnresolvedReference{citation}{[22]} 与 Neri \MBUnresolvedReference{citation}{[13]} 的著作中也可找到这个证明. 关于奇异积分及其应用的综述, 见 A. P. Calderón 的 \emph{Singular integrals}, Bull. Amer. Math. Soc. 72 (1966), 427--465.

\subsubsection{球谐函数}
\label{subsec:ch4-spherical-harmonics}

齐次调和多项式称为实心调和函数, 它在 $S^{n-1}$ 上的限制称为球谐函数. 不同次数的球谐函数在 $L^2(S^{n-1})$ 的内积下正交, 固定次数的球谐函数中也可以选出正交族. 因而 $L^2(S^{n-1})$ 有一个由球谐函数组成的正交基. 当 $n=2$ 时, $L^2(S^1)$ 的球谐基就是 $\{e^{ik\theta}:k\in\mathbb Z\}$.

若 $Y_k$ 是 $k$ 次实心调和函数, 则 $Y_k(x)e^{-\pi|x|^2}$ 是 Fourier 变换的特征函数:
\[
 \widehat{Y_k(x)e^{-\pi|x|^2}}(\xi)
 =(-i)^kY_k(\xi)e^{-\pi|\xi|^2}.
\]
这称为 Bochner--Hecke 公式. 球谐函数在研究奇异积分中的作用可由下面的公式看出. 当 $k\ge1$ 时,
\[
 \left(\operatorname{p.v.}\frac{Y_k(x')}{|x|^n}\right)^{\!\wedge}(\xi)
 =\gamma_k\frac{Y_k(\xi)}{|\xi|^k},
\]
其中 $\gamma_k$ 是只依赖于 $k$ 和 $n$ 的非零常数.

\hypertarget{chapter-04-page-098}{}
对具有偶核 $\Omega\in L^2(S^{n-1})$ 的奇异积分, 可以把 $\Omega$ 展开成球谐函数, 从而给出定理~\ref{thm:ch4-rough-kernel-lp} 的另一个证明, 见 Stein 与 Weiss \MBUnresolvedReference{citation}{[18, Chapter 6]}. 球谐函数的性质和应用见 Stein \MBUnresolvedReference{citation}{[15]}, Stein 与 Weiss \MBUnresolvedReference{citation}{[18]} 以及 Neri \MBUnresolvedReference{citation}{[13]}. 作为调和函数一般理论的一部分, 也可参见 S. Axler, P. Bourdon 与 W. Ramey 的 \emph{Harmonic Function Theory}, Springer-Verlag, New York, 1992.

\subsubsection{算子代数}
\label{subsec:ch4-operator-algebras}

除第~\ref{sec:ch4-operator-algebra}~节研究的算子代数外, Calderón 与 Zygmund 还研究了同类型的其他代数, 见 \emph{Algebras of certain singular integral operators}, Amer. J. Math. 78 (1956), 310--320. 特别地, 令 $\mathcal A_q$ 为所有如下算子的集合:
\[
 Tf(x)=af(x)+\lim_{\varepsilon\to0}\int_{|y|>\varepsilon}
 \frac{\Omega(y')}{|y|^n}f(x-y)\,dy,
\]
其中 $a\in\mathbb C$, $\Omega\in L^q(S^{n-1})$. 若定义
\[
 \|T\|_q=|a|+\|\Omega\|_{L^q(S^{n-1})},
\]
则存在只依赖于 $q$ 的常数 $B_q$, 使得
\[
 \|T_1\circ T_2\|_q\le B_q\|T_1\|_q\|T_2\|_q.
\]
$\mathcal A_q$ 是半单交换 Banach 代数. 关于奇异积分代数的更多内容, 见 A. P. Calderón 的 \emph{Algebras of singular integral operators}, Singular Integrals, Proc. Sympos. Pure Math. X, pp. 18--55, Amer. Math. Soc., Providence, 1967.

\subsubsection{关于旋转方法的进一步结果}
\label{subsec:ch4-more-rotations}

在定理~\ref{thm:ch4-variable-kernel-lq-condition} 的证明中, 若有
\begin{equation}
\label{eq:ch4-mixed-directional-hilbert-bound}
 \left(\int_{\mathbb R^n}
 \left(\int_{S^{n-1}}|H_uf(x)|^{q'}\,d\sigma(u)\right)^{p/q'}dx
 \right)^{1/p}
 \le C\|f\|_p,
\end{equation}
便可由 \eqref{eq:ch4-variable-kernel-holder-bound} 推出 $\|Tf\|_p\le C\|f\|_p$. 在该证明中, 当 $p\ge q'$ 时, \eqref{eq:ch4-mixed-directional-hilbert-bound} 是直接的. 但此式对某些 $p<q'$ 也成立, 从而可以加强定理.

取 $f$ 为球的示性函数, 可以证明 \eqref{eq:ch4-mixed-directional-hilbert-bound} 成立的必要条件是
\[
 \frac np<\frac{n-1}{q'}+1.
\]

\hypertarget{chapter-04-page-099}{}
猜想当 $1<p<\infty$ 且 $p,q'$ 满足这个条件时, \eqref{eq:ch4-mixed-directional-hilbert-bound} 成立. 还猜想在同一指数范围内, 把 $H_u$ 换成相应的极大算子 $H_u^*$ 或方向极大函数 $M_u$, 结论仍成立.

已知结果如下. 当 $n=2$ 时, A. P. Calderón 与 A. Zygmund 在 \emph{On singular integrals with variable kernel}, Appl. Anal. 7 (1978), 221--238 中证明了关于 $H_u$ 的猜想; 当 $n\ge3$ 时, 他们证明了 $1<p<2$ 的情形. M. Cowling 与 G. Mauceri 在 \emph{Inequalities for some maximal functions}, Trans. Amer. Math. Soc. 287 (1985), 431--455 中证明 $M_u$ 也具有相同的指数范围. M. Christ, J. Duoandikoetxea 与 J. L. Rubio de Francia 在 \emph{Maximal operators related to the Radon transform and the Calderón--Zygmund method of rotations}, Duke Math. J. 53 (1986), 189--209 中证明, 对所有这些算子, 猜想在
\[
 1<p<\max\left(2,\frac{n+1}{2}\right)
\]
的范围内成立.

\subsubsection{$L\log L$ 结果}
\label{subsec:ch4-l-log-l-results}

正如第~\ref{sec:ch4-fourier-transform-kernel}~节末尾所指出的, 若偶函数 $\Omega\in L\log L(S^{n-1})$, 则 $\operatorname{p.v.}\Omega(x')/|x|^n$ 的 Fourier 变换有界. 定理~\ref{thm:ch4-rough-kernel-lp} 也可以扩展到偶部属于 $L\log L(S^{n-1})$ 的情形, 见第~\ref{subsec:ch4-references}~节所引 Calderón 与 Zygmund 的第一篇论文.

这些结果是最佳的. 在同一篇论文中, Calderón 与 Zygmund 指出, 若 $\Phi$ 满足
\[
 \frac{\Phi(t)}{t\log t}\longrightarrow0,
 \qquad t\to\infty,
\]
则存在偶函数 $\Omega$, 使 $\Phi(\Omega)\in L^1(S^{n-1})$, 但 $\operatorname{p.v.}\Omega(x')/|x|^n$ 的 Fourier 变换无界.

后来 M. Weiss 与 A. Zygmund 在 \emph{An example in the theory of singular integrals}, Studia Math. 26 (1965), 101--111 中证明, 对每个这样的 $\Phi$, 都存在满足 $\Phi(\Omega)\in L^1(S^{n-1})$ 的偶函数 $\Omega$, 以及对每个 $p>1$ 都属于 $L^p(\mathbb R^n)$ 的连续函数 $f$, 使得对几乎处处的 $x$,
\[
 \limsup_{\varepsilon\to0}
 \int_{|y|>\varepsilon}\frac{\Omega(y')}{|y|^n}f(x-y)\,dy=\infty.
\]

\subsubsection{弱型 $(1,1)$ 不等式}
\label{subsec:ch4-weak-one-one}

命题~\ref{prop:ch4-directional-operator-average} 不能推广到弱型 $(1,1)$ 不等式, 因为 Lorentz 空间 $L^{1,\infty}$ 不是赋范空间, 见第~\ref{subsec:ch2-lorentz-spaces}~节. 所以不能用旋转方法证明 \eqref{eq:ch4-singular-integral-definition} 型奇异积分的弱型 $(1,1)$ 不等式, 甚至不能这样证明 \eqref{eq:ch4-rough-maximal-function} 中粗糙极大函数 $M_\Omega$ 的弱型 $(1,1)$ 不等式. 不过, 已知有如下结果.

\hypertarget{chapter-04-page-100}{}
若 $\Omega\in L\log L(S^{n-1})$, 则 $M_\Omega$ 是弱型 $(1,1)$ 的. $n=2$ 的情形由 M. Christ 在 \emph{Weak type $(1,1)$ bounds for rough operators}, Ann. of Math. 128 (1988), 19--42 中证明; $n\ge2$ 的情形归功于 M. Christ 与 J. L. Rubio de Francia 的 \emph{Weak type $(1,1)$ bounds for rough operators, II}, Invent. Math. 93 (1988), 225--237. 当 $\Omega\in L^1(S^{n-1})$ 时, $M_\Omega$ 是否为弱型 $(1,1)$ 仍属未知.

若 $\Omega\in L\log L(S^{n-1})$, 则 \eqref{eq:ch4-singular-integral-definition} 中的 $T$ 是弱型 $(1,1)$ 的. A. Seeger 在 \emph{Singular integral operators with rough convolution kernels}, J. Amer. Math. Soc. 9 (1996), 95--105 中证明了这一点. M. Christ 与 J. L. Rubio de Francia 在上述论文中, 以及 S. Hofmann 在 \emph{Weak $(1,1)$ boundedness of singular integrals with nonsmooth kernels}, Proc. Amer. Math. Soc. 103 (1988), 260--264 中更早得到部分结果. 由前述反例, 对任意 $\Omega$ 而言这是最佳结果. 不过, 当奇函数 $\Omega\in L^1(S^{n-1})$ 时, $T$ 是否为弱型 $(1,1)$ 仍属未知. P. Sjögren 与 F. Soria 在 \emph{Rough maximal operators and rough singular integral operators applied to integrable radial functions}, Rev. Mat. Iberoamericana 13 (1997), 1--18 中证明, 对任意 $\Omega\in L^1(S^{n-1})$, 将 $T$ 限制到径向函数时是弱型 $(1,1)$ 的.

\subsubsection{分数积分与 Sobolev 空间}
\label{subsec:ch4-fractional-integrals-sobolev}

若 $\varphi\in\mathcal S$, 则
\[
 \widehat{-\Delta\varphi}(\xi)=4\pi^2|\xi|^2\widehat\varphi(\xi).
\]
第~\ref{sec:ch4-operator-algebra}~节用
\[
 \widehat{\Lambda\varphi}(\xi)=2\pi|\xi|\widehat\varphi(\xi)
\]
定义了 $-\Delta$ 的平方根. 更一般地, 可以用
\[
 \widehat{(-\Delta)^{a/2}\varphi}(\xi)
 =(2\pi|\xi|)^a\widehat\varphi(\xi)
\]
定义 Laplace 算子的任意分数次幂. 将它与 \eqref{eq:ch4-fourier-power-kernel} 中 $|x|^{-a}$, $0<a<n$, 的 Fourier 变换比较, 得到所谓分数积分算子 $I_a$, 也称 Riesz 势:
\[
 I_a\varphi(x)=\frac1{\gamma_a}\int_{\mathbb R^n}
 \frac{\varphi(y)}{|x-y|^{n-a}}\,dy,
 \qquad
 \gamma_a=\pi^{n/2-a}
 \frac{\Gamma(a/2)}{\Gamma((n-a)/2)}.
\]
于是, 在分布意义下,
\[
 \widehat{I_a\varphi}(\xi)=|\xi|^{-a}\widehat\varphi(\xi).
\]

\hypertarget{chapter-04-page-101}{}
考察 $I_a$ 在 $L^p$ 上的行为. 齐次性表明范数不等式
\begin{equation}
\label{eq:ch4-fractional-integral-norm}
 \|I_af\|_q\le C\|f\|_p
\end{equation}
只有在
\begin{equation}
\label{eq:ch4-fractional-integral-exponents}
 \frac1q=\frac1p-\frac an
\end{equation}
时才可能成立. 事实上, 对 $1<p<n/a$, 这个条件也是充分的.

\begin{Theorem}
\label{thm:ch4-hardy-littlewood-sobolev}
设 $0<a<n$, $1\le p<n/a$, 并由 \eqref{eq:ch4-fractional-integral-exponents} 定义 $q$. 则 \eqref{eq:ch4-fractional-integral-norm} 成立. 当 $p=1$ 时, $I_a$ 满足弱型 $(1,n/(n-a))$ 不等式
\[
 |\{x\in\mathbb R^n:|I_af(x)|>\lambda\}|
 \le C\left(\frac{\|f\|_1}{\lambda}\right)^{n/(n-a)}.
\]
\end{Theorem}

证明见 Stein \MBUnresolvedReference{citation}{[15, Chapter 5]}. 另一个证明使用 L. Hedberg 在 \emph{On certain convolution inequalities}, Proc. Amer. Math. Soc. 36 (1972), 505--510 中给出的不等式
\begin{equation}
\label{eq:ch4-hedberg-inequality}
 I_af(x)\le C_a\|f\|_p^{ap/n}Mf(x)^{1-ap/n},
 \qquad 1\le p<n/a,
\end{equation}
其中 $M$ 是 Hardy--Littlewood 极大函数. 为证明它, 将定义 $I_a$ 的积分分成 $|x-y|\le A$ 与 $|x-y|>A$ 两部分. 对第一部分应用命题~\ref{prop:ch2-radial-majorization}, 对第二部分应用 Hölder 不等式, 再选取 $A$ 使两项大小相同. 定理~\ref{thm:ch4-hardy-littlewood-sobolev} 随即由 \eqref{eq:ch4-hedberg-inequality} 和 $M$ 的有界性推出.

与分数积分算子密切相关的是分数极大函数
\[
 M_af(x)=\sup_{r>0}\frac1{|B_r|^{1-a/n}}
 \int_{B_r}|f(x-y)|\,dy,
 \qquad 0<a<n.
\]
$M_a$ 满足与 $I_a$ 相同的范数不等式. 弱型 $(1,n/(n-a))$ 不等式可用覆盖引理论证证明, 见第~\ref{subsec:ch2-covering-lemmas}~节. 又由 Hölder 不等式, $M_a$ 是强型 $(n/a,\infty)$ 的. 强型 $(p,q)$ 不等式于是由插值得到.

对任意正函数 $f$, 容易看出 $M_af(x)\le CI_af(x)$. 反向的逐点不等式一般不成立, 但两者的范数可比较.

\hypertarget{chapter-04-page-102}{}
\begin{Theorem}
\label{thm:ch4-fractional-integral-maximal-comparison}
若 $1<p<\infty$, $0<a<n$, 则存在常数 $C_{a,p}$, 使得
\[
 \|I_af\|_p\le C_{a,p}\|M_af\|_p.
\]
\end{Theorem}

这个结果归功于 B. Muckenhoupt 与 R. Wheeden 的 \emph{Weighted norm inequalities for fractional integrals}, Trans. Amer. Math. Soc. 192 (1974), 261--274. 另见 D. R. Adams 的 \emph{A note on Riesz potentials}, Duke Math. J. 42 (1975), 765--778, 以及 D. R. Adams 与 L. Hedberg 的 \emph{Function Spaces and Potential Theory}, Springer-Verlag, Berlin, 1996.

定理~\ref{thm:ch4-hardy-littlewood-sobolev} 的一个重要应用是证明 Sobolev 嵌入定理. 对 $1\le p<\infty$ 和整数 $k\ge1$, 定义 Sobolev 空间 $L_k^p(\mathbb R^n)$ 为所有这样的函数 $f$ 的空间: $f$ 及其不超过 $k$ 阶的所有分布导数都属于 $L^p$. 更准确地说, $L_k^p$ 是 $C_c^\infty$ 在范数
\[
 \|f\|_{L_k^p}=\sum_{|\alpha|\le k}\|D^\alpha f\|_p
\]
下的闭包.

\begin{Theorem}
\label{thm:ch4-sobolev-embedding}
固定 $1\le p<\infty$ 和整数 $k\ge1$. 则:
\begin{enumerate}
\item 若 $p<n/k$, 则对所有 $p\le r\le q$, $L_k^p(\mathbb R^n)\subset L^r(\mathbb R^n)$, 其中 $1/q=1/p-k/n$;
\item 若 $p=n/k$, 则对所有 $p\le r<\infty$, $L_k^p(\mathbb R^n)\subset L^r(\mathbb R^n)$;
\item 若 $p>n/k$ 且 $f\in L_k^p(\mathbb R^n)$, 则 $f$ 在一个零测集外与连续函数相等.
\end{enumerate}
\end{Theorem}

这个结果归功于 S. L. Sobolev 的 \emph{On a theorem in functional analysis}, Mat. Sb. 46 (1938), 471--497; 英译文见 Amer. Math. Soc. Transl. ser. 2, 34 (1963), 39--68. 本小节的全部结果可见 Stein \MBUnresolvedReference{citation}{[15, Chapter 5]}, 上引 Adams 与 Hedberg 的著作, 以及 R. Adams 的 \emph{Sobolev Spaces}, Academic Press, New York, 1975 和 W. Ziemer 的 \emph{Weakly Differentiable Functions}, Springer-Verlag, New York, 1989.
